% Generated draft: requires mathematical and source review.
\documentclass[11pt]{article}
\usepackage[margin=25mm]{geometry}
\usepackage{fontspec}
\setmainfont{Latin Modern Roman}
\usepackage{amsmath,amssymb,mathtools,unicode-math}
\setmathfont{Latin Modern Math}
\usepackage{enumitem,needspace,xurl,hyperref}
\setlength{\parindent}{0pt}
\setlength{\parskip}{5pt}
\setlength{\emergencystretch}{3em}
\newcommand{\N}{\mathbb{N}}
\newcommand{\Z}{\mathbb{Z}}
\newcommand{\Q}{\mathbb{Q}}
\newcommand{\R}{\mathbb{R}}
\newcommand{\C}{\mathbb{C}}
\newcommand{\F}{\mathbb{F}}
\newcommand{\A}{\mathbb{A}}
\newcommand{\PP}{\mathbb{P}}
\newcommand{\E}{\mathbb{E}}
\newcommand{\eps}{\varepsilon}
\newcommand{\dd}{\,\mathrm{d}}
\begin{document}
\section*{UnsolvedMath ID: 30006289}
\textbf{Source problem code:} catalogue entry
\subsection{Definitions and mathematical statement}
The supplied excerpt concerns an Euler system whose equations, spatial domain, constitutive assumptions, initial data, and boundary conditions are not included. It also does not define the admissible numerical schemes, convergence topology, dissipative solution class, entropy-production comparison, or computability requirement. Consequently, the following is a conditional statement schema requiring these specifications before it becomes a precise mathematical conjecture.

Let \(\Sigma\) denote the intended class of numerical schemes and \(\mathcal X\) the intended space of candidate solutions. For \(\mathsf S\in\Sigma\) and \(U\in\mathcal X\), use the following placeholders:
\begin{itemize}
\item \(W(\mathsf S,U)\): the approximations produced by \(\mathsf S\) converge, in the intended topology and refinement regime, to \(U\), and \(U\) is a weak solution of the intended Euler system.
\item \(O(\mathsf S)\): the scheme \(\mathsf S\) is oscillatory in the intended sense.
\item \(M(\mathsf S,U)\): \(U\) is a maximal computable dissipative solution in the intended comparison class associated with the assertion about \(\mathsf S\).
\end{itemize}
These predicates record the source terminology; they do not supply its missing mathematical definitions. In particular, the notation \(M(\mathsf S,U)\) leaves unresolved whether computability or the comparison class depends on the scheme.

The two assertions in the supplied conjecture have the schematic form
\[
\forall\mathsf S\in\Sigma\;\forall U\in\mathcal X,
\qquad W(\mathsf S,U)\Longrightarrow M(\mathsf S,U),
\]
and
\[
\forall\mathsf S\in\Sigma,
\qquad O(\mathsf S)\Longrightarrow
\neg\exists U\in\mathcal X\;M(\mathsf S,U).
\]
The second implication concerns the same numerical scheme as the first. Its conclusion reproduces the source's phrase ``maximal (computable) dissipative solution''; whether the parentheses weaken or merely qualify the computability requirement is unresolved. The source discusses maximality in connection with entropy-production rate, but supplies neither its ordering nor the comparison class, and does not distinguish maximality from attainment of a maximum.

The surrounding excerpt reports the following context, with hypotheses unavailable in the supplied record: dissipative solutions coincide with an existing strong solution throughout its lifespan, citing [1]; approximations converging weakly but not strongly have a limit that is not a weak solution and have nonzero defects; \(K\)-convergence gives strong convergence of empirical means of numerical solutions to observable quantities of a dissipative solution, citing [4] and [5]; and selecting a dissipative solution maximal with respect to entropy-production rate makes the defects vanish and yields a weak solution. It also reports that standard structure-preserving numerical methods applied to oscillatory solutions produce dissipative solutions with nonzero defects. These are supplied contextual claims, not additional conclusions established here. In particular, the relationship between this reported defect-vanishing selection and the conjectured nonexistence requires the missing definitions and hypotheses.
\subsection{Short English statement}
The conjecture links a numerical scheme's behavior to maximal computable dissipative solutions of the Euler equations: convergence to a weak solution should produce such a maximal solution, while oscillation of the same scheme should imply nonexistence. The supplied excerpt omits essential definitions and hypotheses, so only a conditional formulation is possible. Its open status is supplied curation, and this statement-only draft requires human mathematical review.
\subsection{Sources}
\begin{itemize}
\item Local UnsolvedMath registry: supplied record with permanent ID \(30006289\), title \emph{Maximal Computable Dissipative Solutions of Euler Equations}, source problem code \texttt{catalogue entry}, in the Oberwolfach Reports Open Problems collection. The recorded proposed year is 2025; no proposer is identified.
\item Supplied report citation: \emph{Computational Multiscale Methods} (2025), Oberwolfach Reports, pp.\ 1079--1162, DOI \texttt{10.4171/owr/2025/22}. Source URL and literature-source link labelled ``original OWR report'': \url{https://doi.org/10.4171/owr/2025/22}. This is report-level provenance; no conjecture-specific page is supplied.
\item Supplied literature-source attribution: Feireisl--Lukacova-Medvidova--Yu, ``oscillatory approximations'': \url{https://arxiv.org/abs/2505.02070}. The record provides no fuller bibliographic description or explicit correspondence with the excerpt's numbered references.
\item The original excerpt attributes the conjecture to ``our recent work [6]'' and uses references [1], [4], and [5] for contextual claims. Their bibliographic entries and the authors represented by ``our'' are not supplied.
\item Literature assessment dated exactly \texttt{2026-08-22}: the supplied curation labels the problem open and describes the assertion as a recent conjecture, stating that existing \(K\)-convergence and entropy-selection results do not prove the two-part assertion for arbitrary Euler schemes. This is supplied dated curation, not an independently verified literature assessment; the record requests independent verification of time-sensitive status and cited claims.
\item The registry's \texttt{corrected\_verified} statement label is accompanied by a qualification that only clear formula delimiters, notation, grammar, and typography were repaired while mathematical content and scope were preserved. It is not treated here as independent mathematical verification.
\end{itemize}
\end{document}
